\chapter{Проблема}

\section{Как сегодня устроена агентная разработка}

Почти все инструменты, в которых агенты пишут код, собраны вокруг одной схемы. Агент запускается там, где лежит код. Изоляция параллельной работы достигается копиями рабочего дерева: git worktree на локальной машине или отдельная виртуальная машина в облаке. Параллелизм достигается умножением инстансов: десять задач, десять агентов, десять копий среды.

Для веб-сервисов и бэкендов эта схема работает хорошо. Среда там дешёвая: контейнер с компилятором и зависимостями поднимается за секунды, и облачные песочницы на каждого агента стоят копейки. Ровно поэтому индустрия оптимизирует именно эту нишу, и ровно поэтому она не смотрит на другие.

Схема наследует свою главную предпосылку у git. Git создавался для людей и предполагает, что у каждого разработчика есть собственная копия среды. Для человека это разумно: людей мало, их машины сильные, и каждый работает в своём темпе. Для агентов та же предпосылка превращается в три проблемы.

\section{Три проблемы}

\subsection{Среда дорогая или невоспроизводимая}

Сборка iOS-приложения требует macOS и Xcode конкретной версии. Подпись требует ключей, которые нельзя раздавать. Тесты на устройствах требуют самих устройств. Обучение или инференс модели требует GPU, которых нет в дешёвой песочнице. Работа с данными компании требует, чтобы данные не покидали её периметр. Большой монорепозиторий требует часов на прогрев кешей сборки.

Во всех этих случаях воспроизвести среду на каждого агента либо очень дорого, либо невозможно. Схема с копией среды на инстанс здесь просто не работает, и команды остаются с одним агентом на одной машине.

\subsection{Лимиты провайдера на одного владельца}

Подписка на модель выдаёт квоту одному владельцу: столько-то токенов в окно времени. Десять агентов под одной подпиской делят одну квоту, и рой упирается в лимит раньше, чем в вычислительную мощность. Обойти это можно только большим числом подписок, и здесь начинается третья проблема.

\subsection{Чужие агенты на одной машине означают чужие учётные данные на ней}

Пусть над проектом работают пять человек, у каждого своя подписка. Чтобы их агенты работали на одной машине с дорогой средой, каждый должен передать этой машине свои учётные данные. Условия провайдеров это запрещают, а здравый смысл не советует: машина, на которой лежат чужие ключи, становится целью. Значит, объединить подписки нескольких людей на одной среде сегодня нельзя, не нарушив либо правила, либо безопасность.

\section{Неравенство}

Из трёх проблем следует неравенство, которое и есть настоящий предмет проекта.

Большие рои агентов, координируемые для работы над одним проектом, доступны сегодня вендорам моделей, владельцам крупных вычислительных ресурсов и корпорациям. У них есть и квота, и инфраструктура, и возможность поставить среду рядом с агентами. Открытые проекты активно используют агентов, но спроектированы под координацию через людей: issue, pull request, ревью мейнтейнера. Агент в открытом проекте это помощник одного человека, а не участник роя.

Мы предполагаем, что координация на уровне автономного роя эффективнее координации через людей: человек читает pull request часами, рой обменивается назначениями и вердиктами за секунды. Это гипотеза, и глава 11 описывает эксперимент, который её проверяет; в работе она нигде не утверждается как факт. Разрыв между теми, кто может собрать рой, и всеми остальными в любом случае растёт вместе с возможностями моделей.

Проект есть ответ на это неравенство: протокол, который позволяет незнакомым друг с другом людям приносить в общий проект не патчи, а ёмкость своих подписок и моделей, и делает это безопасным для проекта и для них самих.

\section{Что мы утверждаем и чего не утверждаем}

Важно с самого начала провести границу честно, потому что от неё зависит, как оценивать результат.

Мы не утверждаем, что протокол делает агентов умнее. Качество декомпозиции задач, вероятность того, что ревьювер заметит дефект, и сила тестового набора остаются свойствами моделей и проекта. Протокол ни одну из этих трёх величин не улучшает.

Мы утверждаем другое, и формулировка вклада в третьей редакции работы звучит так: мы рассматриваем совместное изменение проекта агентами независимых владельцев в домене, где у принятия результата нет механического оракула; рабочая нода может предлагать и производить изменения, но не может расширить область задачи, изменить условия её приёмки или продвинуть основную ветку; авторизация, исполнение и принятие связаны проверяемыми записями, относящимися к конкретным версиям контракта задачи и состояния проекта. Иначе говоря, никакой агент не получает больше, чем ему выдано, и ничего не попадает в основную ветку без независимой проверки ровно того состояния, которое попадает. Протокол делает суждение агентов независимым, случайным и подотчётным: ревьювер выбран так, что его нельзя подобрать, и каждое действие оставляет след, по которому его можно воспроизвести. Протокол делает возможным то, что сегодня доступно немногим: рой из ёмкости независимых людей.

Ценность здесь это доступ и контроль. Эффективность нужно измерять, а не декларировать: принятые задачи на миллион токенов, время от постановки до интеграции, доля откатов, часы человеческого ревью на принятую задачу. Эти метрики встроены в систему с первого дня, и глава 12 объясняет, как они собираются.

\section{Сценарий}

Чтобы дальнейший текст не был абстрактным, зафиксируем сценарий, к которому книга будет возвращаться.

\begin{exbox}[Сценарий: приложение с нативным ядром]
Проект: iOS-приложение с ядром на Rust. Среда исполнения: один Mac Studio с Xcode, ключами подписи и подключёнными устройствами, стоящий у кор-команды. Кор-команда: два человека, у каждого своя кор-нода. Контрибьюторы: пять человек в разных городах, у троих подписки на разные модели, у одного ключ API, у одного локальная модель на собственном Mac. Никто из них не имеет доступа к Mac Studio, ключам подписи или устройствам.

Кор-команда ставит цель: перенести модуль синхронизации на новый протокол. Оркестратор, одна из кор-нод, раскладывает цель на двадцать задач с областями изменений и критериями приёмки. Каждую задачу подхватывает агент одного из контрибьюторов, работает в изолированном каталоге на Mac Studio через удалённый shell, представляет коммит. Случайно выбранный агент другого контрибьютора делает ревью. Тесты гоняет третий. Среда сливает результат в основную ветку только тогда, когда журнал подтвердил, что ревью и тесты относятся ровно к тому, что сливается. Если что-то идёт не так, любая кор-нода останавливает всё одной подписью.
\end{exbox}

Всё, что дальше, есть объяснение того, как этот сценарий работает и почему его нельзя сломать, не сломав явных допущений.

\section{Контекст: что уже сделали другие}

Объединение чужих агентов на общем проекте само по себе уже не ново, и это важно понимать до чтения остального. В формальной математике открытая платформа Prove2Me позволяет любому владельцу агента вносить доказательства в общую миссию: граф формулировок задаёт фронт работы, формулировки отделены от доказательств, описание каждого узла делает результаты находимыми и переиспользуемыми. На ней в сентябре 2026 года была завершена формализация Великой теоремы Ферма силами десятков агентов, а в отдельном эксперименте три личные подписки на модель за три дня совместно формализовали крупный результат. Для обычной разработки ближайшая работа это команда из шестнадцати параллельных агентов, построившая компилятор Си: агенты в отдельных контейнерах, общий репозиторий, блокировки задач, самостоятельная интеграция.

Чем мы отличаемся. В формальной математике у принятия есть оракул, ядро проверки доказательств, и вся координация стоит на нём. У кода оракула нет, и принятие приходится конструировать: отсюда независимое ревью, кандидат, аттестации и кворум. В эксперименте с компилятором все агенты одного владельца и доверены; у нас нет. Наш контракт задачи и граф зависимостей есть перенос разделения формулировок и доказательств на код, а всё остальное есть ответ на отсутствие оракула и на недоверие к участникам.
